\section{Introduction}\label{sec:introduction}

\subsection{History and scope}\label{hist:history}

Penrose's 1973 argument connected gravitational collapse and cosmic censorship
with a lower bound for the mass in terms of black-hole area
\cite{penrose1973}.  In the asymptotically flat, time-symmetric setting,
Huisken and Ilmanen proved the three-dimensional inequality for the area of
each connected outermost minimal component by weak inverse mean curvature
flow \cite{huiskenilmanen2001}.  Bray's conformal flow proved the bound for
the total area of a possibly disconnected outermost minimal boundary
\cite{bray2001}; Bray and Lee extended the Riemannian inequality to dimensions
less than eight \cite{braylee2009}.  These are scalar-curvature theorems for
Riemannian metrics.  Recovering an arbitrary original second fundamental
form in an equality statement requires further arguments.

For cosmological constant $-3$, the Schwarzschild--anti-de Sitter horizon
relation is $M=(r_h+r_h^3)/2$.  It leads to the asymptotically hyperbolic
Penrose bound, whose mass convention comes from the metric mass integrals
of Wang and Chru\'sciel--Herzlich \cite{wang2001,chruscielherzlich2003}.
Earlier Riemannian results include the graphical estimates of Dahl,
Gicquaud and Sakovich \cite{dahlgicquaudsakovich2013}, the sharp Inequality
of de Lima and Gir\~ao for balanced hyperbolic graphs with scalar curvature
at least $-6$ and a minimal horizon in a horizontal totally geodesic slice
met orthogonally, whose projection is star-shaped and mean convex
\cite[Theorem~1.2]{delimagirao2016},
and Ambrozio's Theorem for sufficiently small perturbations of a fixed
positive-mass Schwarzschild--anti-de Sitter metric, with minimal boundary
and scalar curvature at least $-6$ \cite{ambrozio2015}.
Lee and Neves treat a different asymptotically locally hyperbolic regime,
using a nonpositive supremum of the mass aspect and a boundary component
with the genus of conformal infinity \cite{leeneves2015}.
For static systems, Borghini, Fogagnolo and Pinamonti characterize
finite domains attaining equality in the substatic Heintze--Karcher
inequality as warped products
\cite[Theorem~1.1]{borghinifogagnolopinamonti2024}.
We use this finite-domain result after proving that the equality data
force the static deficit to vanish, and then prove the global extension.
The flow issue at infinity is substantive: Neves constructed an
asymptotically hyperbolic inverse mean curvature flow lacking the convergence
needed for the direct Hawking-mass argument \cite{neves2010}.

There are also perturbative results involving the second fundamental form.
Khuri and Kopi\'nski prove the inequality for small maximal vacuum conformal
perturbations of Schwarzschild--anti-de Sitter data under a quantitative
comparison between a weighted bulk norm and a boundary-related norm of the
seed tensor \cite{khurikopinski2023}.
The earlier method manuscript \cite{internalads} treats a fixed seed and a
fixed local maximal vacuum conformal branch; its quadratic-null directions
require a fourth-order calculation.  The asymptotically flat manuscript
\cite{internalflat} develops coupled deformation and original-data rigidity
methods at zero cosmological constant.  We use these as methodological
antecedents and supply the arguments needed at the present hypotheses.

\subsection{Significance and technical contributions}\label{hist:significance}

We establish the three-dimensional nonmaximal asymptotically anti-de Sitter
spacetime Penrose Inequality for the precise class specified below.
Its numerical assertion uses the full enclosing-area infimum
of a possibly disconnected weakly future outer trapped
boundary and the Lorentz norm of the four \emph{metric} fluxes.
In the connected, outermost, outer area-minimizing future marginally
outer trapped subclass, its equality assertion recovers the original pair
$(g,K)$ as a smooth proper spacelike hypersurface in the smooth maximal
extension of Schwarzschild--anti-de Sitter spacetime of parameter exactly
$M=m_{\mathrm{AH}}$, with a full future horizon section (the bifurcation
sphere is allowed) and an end reaching conformal infinity.
This resolves the asymptotically anti-de Sitter spacetime Penrose Conjecture
in this initial-data class, including its Riemannian specialization and
original-data rigidity on the stated connected horizon subclass.
The mass match concerns the designated initial
slice; spatial chart covariance supplies no invariance under arbitrary
changes of that slice.  The assertions are initial-data statements and
make no claim to prove cosmic censorship or an evolution theorem.

Three parts of the construction have uses beyond the model calculation.
First, the coupled deformation combines weak trapping barriers with
estimates that distinguish the constants required uniformly in the large
profile parameter from those needed only after that parameter is fixed.
A scalar regularity argument for a measurable rank-one axial coefficient
and a bounded-input solution operator provide the analytic ingredients
for the degree construction.  Second, a primitive correcting the spatially
varying lapse cancels mixed derivatives in the hyperbolic deformation;
the resulting weighted divergence estimate controls all four mass fluxes
and retains the full enclosing-area comparison.  Third, equality is tested
on the original constraints.  The first stationarity argument determines
the horizon surface gravity, and a causal adjoint with analysis across its
null set produces the complete static quotient.  Weighted improvement of
its end evaluates the static Heintze--Karcher deficit in terms of the
original mass. The equality data force its nonnegative interior defect
to vanish. A finite horizon collar then attains exact equality, and the
static equations give global Schwarzschild--anti-de Sitter identification.
The remaining reconstruction recovers both original tensors and the proper
hypersurface embedding, including its future horizon attachment and
conformal infinity.


\subsection{The results}
The mass in this paper is the hyperbolic metric four-flux mass.
The area is the infimum over full enclosing cuts, including all intrinsic
boundary components. Precise definitions appear in Section~\ref{sec:setup}.

\begin{theorem}[Nonmaximal anti-de Sitter Penrose Inequality]\label{thm:main-numerical}
Let $(\Omega,g,K)$ satisfy Definition~\ref{def:setup-data}, the same
future boundary condition $H+\tr_{TS}K\le0$ on every component, and
the future-timelike metric mass condition in Definition~\ref{def:setup-mass}.
Then
\begin{equation}\label{eq:main-penrose}
 m_{\AH}(g)\ge
 \sqrt{\frac{A_{\min}(S)}{16\pi}}
 \left(1+\frac{A_{\min}(S)}{4\pi}\right)
 =\frac{r_A+r_A^3}{2}.
\end{equation}
The coefficient is sharp. No maximality assumption, restriction on the
compact topology or boundary genera, or outermostness hypothesis is required.
\end{theorem}

\begin{theorem}[Original-data equality]\label{thm:main-rigidity}
Assume in addition the connected horizon hypotheses of
Definition~\ref{def:setup-horizon}. Equality in Equation~\ref{eq:main-penrose}
holds if and only if the original pair $(g,K)$ is realized by a smooth
proper spacelike embedding of $\Omega$, with boundary, into the smooth
maximally extended Schwarzschild--anti-de Sitter spacetime of parameter
exactly $M=m_{\AH}(g)$. Its interior lies in one black-hole exterior;
its boundary maps diffeomorphically onto a full smooth spacelike
cross-section of that exterior's future event horizon, with the
bifurcation sphere also allowed; and its unique end reaches that
exterior's conformal infinity. The embedding induces both $g$ and $K$
with the future-normal convention of Equation~\ref{eq:setup-expansion}.
\end{theorem}

The mass match in Theorem~\ref{thm:main-rigidity} is part of the statement.
It is not a consequence of being an arbitrary asymptotic time slice of
a Schwarzschild--AdS spacetime. The theorem permits all smooth slices
satisfying its intrinsic hypotheses and this computed mass match.
The numerical theorem also allows nonoutermost boundaries; its equality
configurations outside the connected horizon subclass are not classified here.

\subsection{Proof architecture}
The numerical argument first establishes a fixed-chart time-component
bound without assuming that its comparison mass is timelike. In the
one-sign case, a conserved auxiliary stress has a flux that approaches
the minimum enclosing area. Removing this stress at infinity changes
the asymptotic model and produces the cubic area term. The resulting
asymptotically flat data are handled by a coupled elliptic deformation
and the Riemannian Penrose Inequality. A separate weighted deformation
on the hyperbolic end removes a general original second fundamental form.
Spatial Lorentz covariance then gives Equation~\ref{eq:main-penrose}.

The two coupled constructions require global existence, not just
ellipticity of their scalar equations. We prove the necessary bounds,
the boundary version of a rank-one measurable-coefficient gradient
estimate, the bounded-input scalar solve, and degree invariance on the
moving admissible regions. Every final small error uses an explicit
parameter order. Constants used only for fixed-parameter compactness
are kept separate from those that must be polynomial in the parameter.

Section~\ref{sec:fl-system-floor} states the flat estimate and constructs
its equations and lower-bound barriers. Section~\ref{reg-section}
supplies the local regularity estimate used in both deformations, and
Section~\ref{sec:fc-existence} proves existence and completes the flat
mass comparison. Section~\ref{sec:ah-deformation} then adapts these tools
to the hyperbolic end and proves the general numerical inequality.

For equality, variations of the original data produce a causal
stationary field and identify the boundary conditions, including the
surface gravity. Twist estimates and exclusion of an interior null set
give a static quotient. Its normalized end, horizon area and surface
gravity make the static Heintze--Karcher deficit vanish. A bounded
elliptic exhaustion and the finite-domain equality theorem of
Borghini--Fogagnolo--Pinamonti
\cite[Theorem~1.1]{borghinifogagnolopinamonti2024} identify a warped
horizon collar; the static equations extend this identification over
the complete exterior. Finally we recover the original
spacelike graph, including its future horizon boundary and its second
fundamental form. Neither equality of a deformation metric nor an
assumed spacetime realization replaces this last step.

Figure~\ref{fig:proof-architecture} records these dependencies.
\begin{figure}[tbp]
\centering
\begin{tikzpicture}[
  archbox/.style={draw=black!65, rounded corners=2pt, line width=.45pt,
    text width=2.94cm, minimum height=1.42cm, inner sep=3.5pt,
    align=center, font=\footnotesize},
  archstep/.style={archbox, fill=black!2},
  archresult/.style={archbox, fill=blue!5},
  archarrow/.style={-{Stealth[length=1.7mm,width=1.2mm]},
    draw=black!75, line width=.55pt},
  archheading/.style={anchor=west, font=\small\bfseries},
  archlabel/.style={font=\scriptsize, fill=white, inner sep=2pt}
]
\node[archheading] at (-1.60,1.26)
  {(a) Numerical inequality};
\node[archstep] (archstress) at (0,0)
  {Auxiliary stress\\ Conserved; flux saturation\\
   Lemmas~\ref{lem:str-flux-saturation}, \ref{lem:str-stress-identities}};
\node[archstep] (archbend) at (3.85,0)
  {Bend to a flat end\\ Cubic mass correction\\
   Proposition~\ref{prop:str-af-bending}};
\node[archstep] (archflat) at (7.70,0)
  {AF coupled deformation\\
   $R_{\widehat g}\ge0$\\
   Lemma~\ref{lem:fc-comparison-metric}};
\node[archstep] (archpenrose) at (11.55,0)
  {Riemannian Penrose\\ Flat mass estimate\\
   Theorem~\ref{thm:fc-mass-completion}};
\draw[archarrow] (archstress) -- (archbend);
\draw[archarrow] (archbend) -- (archflat);
\draw[archarrow] (archflat) -- (archpenrose);

\node[archstep] (archah) at (0,-2.50)
  {AH deformation\\
   $(\widehat g,0)$\\ Mass and area control\\
   Proposition~\ref{prop:ah-existence}};
\node[archresult] (archonesign) at (3.85,-2.50)
  {One-sign bound\\ Applied to $(\widehat g,0)$\\
   Theorem~\ref{thm:str-one-sign}};
\node[archresult] (archcomponent) at (7.70,-2.50)
  {Fixed-chart bound for $g$\\
   $p_0(g)\ge (r_A+r_A^3)/2$\\
   Theorem~\ref{thm:ah-component}};
\node[archresult] (archlorentz) at (11.55,-2.50)
  {Spatial Lorentz\\ Balancing\\
   $m_{\AH}\ge (r_A+r_A^3)/2$\\
   Theorem~\ref{thm:ah-invariant}};
\draw[archarrow] (archpenrose.south) -- (11.55,-1.22)
  -- node[archlabel,above] {ordered limits of inequalities}
  (3.85,-1.22) -- (archonesign.north);
\draw[archarrow] (archah) -- (archonesign);
\draw[archarrow] (archonesign) -- (archcomponent);
\draw[archarrow] (archcomponent) -- (archlorentz);

\node[archheading] at (-1.60,-3.80)
  {(b) Equality on the connected horizon subclass};
\node[archstep] (archvariations) at (0,-5.02)
  {Original-data variations\\ At equality in $(g,K)$\\
   Section~\ref{eq-section}};
\node[archstep] (archadjoint) at (3.85,-5.02)
  {Causal adjoint\\ $(u,X)$; horizon data\\
   Theorem~\ref{eq-causal-killing}};
\node[archstep] (archstatic) at (7.70,-5.02)
  {No interior null set\\ Regular static quotient\\
   Theorem~\ref{geo-static-quotient}};
\node[archstep] (archend) at (11.55,-5.02)
  {Weighted end estimates\\ Normalized smooth end\\
   Proposition~\ref{st-end-regularity}};
\draw[archarrow] (archvariations) -- (archadjoint);
\draw[archarrow] (archadjoint) -- (archstatic);
\draw[archarrow] (archstatic) -- (archend);

\node[archstep] (archhorizon) at (3.85,-7.77)
  {Horizon area and gravity\\
   $|S|_h=4\pi r_h^2$\\
   $\kappa=\frac{1+3r_h^2}{2r_h}$\\
   Proposition~\ref{st-base-equality}};
\node[archresult] (archhw) at (7.70,-7.77)
  {Vanishing static deficit\\ Global static model\\
   Theorem~\ref{hw-static-rigidity}};
\node[archresult] (archgraph) at (11.55,-7.77)
  {Original graph\\
   $(g,K)$, $M=m_{\AH}$\\ Full future horizon\\ Conformal infinity\\
   Proposition~\ref{emb-global}};
\draw[archarrow] (archstatic.south) -- (7.70,-6.16)
  -- (3.85,-6.16) -- (archhorizon.north);
\draw[archarrow] (archend.south) -- (11.55,-6.44)
  -- (7.70,-6.44) -- (archhw.north);
\draw[archarrow] (archhorizon) -- (archhw);
\draw[archarrow] (archhw) -- (archgraph);
\end{tikzpicture}
\caption{Dependencies of the two arguments. AF means asymptotically flat
and AH means asymptotically hyperbolic. Arrows indicate input to the
next step. In (a), the one-sign theorem is first proved by the upper
route and then applied to the AH comparison data; their mass covector
need not be timelike. Only the original mass covector is balanced.
In (b), equality and the numerical inequality give variations of the
original data. The normalized end and the inherited horizon area and
surface gravity jointly supply the static rigidity hypotheses.
Numerical area comparisons use full enclosing cuts.
The equality route uses no limit of the deformation metrics.}
\label{fig:proof-architecture}
\end{figure}


The local conformal AdS and asymptotically flat rigidity manuscripts
\cite{internalads,internalflat} motivate parts of the construction.
No theorem from either is assumed: the reused arguments and all changes
required by the present hypotheses are proved below. The public inputs
are invoked with their local geometry, regularity, asymptotic and
boundary hypotheses checked where they are used.
